\section{Lean, literature, and lessons for automation}\label{sec:lean-lit}

\subsection{Lean status}

The five arches are proved in Lean~4 with Mathlib in \lp{}, with no axioms beyond \texttt{propext}, \texttt{Classical.choice} and \texttt{Quot.sound}. This describes the development: the files contain no \texttt{\#print axioms} command, and we do not record a separate check. The generated statements are formalised as corollaries of the arches in the Lake package \texttt{lean/bridges} of the repository of this edition (\url{https://github.com/Commutator-IO/grothendieck-project}), which requires lean-pool at a fixed revision. The package has the following files:
\begin{itemize}
\item \lean{PreorderPresheaves.lean} (\lean{exists_preorder_moritaEquivalentCat_iff_isThin}, after Grothendieck's exercise in folder 161-3, p.~26);
\item \lean{Monoids.lean} (\lean{isBooleanTopos_singleObj_iff}, \lean{isDeMorganTopos_singleObj_iff}, \lean{isDeMorganTopos_singleObj_of_monoidWithZero}, \lean{isDeMorganTopos_singleObj_of_commMonoid}, \lean{isDeMorganTopos_singleObj_freeMonoid_unit}, \lean{not_isDeMorganTopos_freeMonoid_two}, \lean{isBooleanTopos_singleObj_of_group}); \lean{Preorders.lean} (\lean{isDeMorganTopos_functor_preorder_iff}, \lean{extremallyDisconnected_iff_confluent}, \lean{moritaEquivalentCat_preorder_iff}); \lean{Combinatorics.lean} (\lean{isDeMorganTopos_presheaf_of_hasInitial}, \lean{not_isDeMorganTopos_graphs}, \lean{graphs_pseudocompl_counterexample}, \lean{not_isDeMorganTopos_simplicialSets}, \lean{isDeMorganTopos_augmentedSimplicialSets}, \lean{isDeMorganTopos_presheaf_fintypeCatSkeleton}, \lean{isDeMorganTopos_treeTopos}, \lean{not_isBooleanTopos_treeTopos}); \lean{EssentialPoints.lean} (\lean{nonempty_essentialPoint_iso_iff}, \lean{exists_iso_essentialPoint}, \lean{essentialPointEndMulEquivOfGroup}, \lean{karoubiRel_iff_rank_eq}, \lean{essentialPointClassesEquivFin}).
\item \lean{QuiversMorita.lean} (Morita and essential points of quiver representations: \lean{moritaEquivalentCat_paths_iff}, \lean{nonempty_representations_equiv_iff}, \lean{essentialPointPathsEquiv}) and \lean{QuiversCyclic.lean} (the $2^n+1$ topologies of the cyclic quiver: \lean{card_grothendieckTopology}).
\item \lean{QuiversDeMorgan.lean} (De Morgan and Boolean for quiver representations: \lean{rightOreCondition_paths_iff}, \lean{isDeMorganTopos_paths_iff}, \lean{isBooleanTopos_paths_iff}; it contains no Morita or essential-point statement), \lean{Corners.lean} (transfer to corners), \lean{KaroubiMonoid.lean} (the $n$ classes for $T_n$), and, for the calibration log, \lean{FiniteMonoids.lean} (single minimal right ideal) and \lean{NatTopologies.lean} (the three topologies on $\B\mathbb N$, \lean{nat_card_grothendieckTopology_natCat}, with the sheaves of the middle one, \lean{isSheaf_finalTopology_iff}).
\end{itemize}
Not formalised: the analysis chapter (Mathlib has no topos of sheaves on a space with its negation), the subspace lattices beyond the De Morgan and Morita statements, the counts for $PT_n$, $I_n$, $O_n$ and $(\mathbb Z/n,\cdot)$, all the counts of subtopoi of monoids, and the description of the sheaves of the cyclic quiver. The machine enumerations for $(\mathbb Z/n,\cdot)$ (\S\ref{sec:linear}) are not part of the Lean development, and for $n>30$ they go through the classification of idempotent ideals, so they do not check the topologies independently.

In Mathlib, \texttt{SingleObj M} composes as $f\gg g = g\cdot f$; with this convention the Ore condition on \texttt{SingleObj M} reads $\forall a\,b,\ \exists c\,d,\ a\cdot c = b\cdot d$ (\lean{rightOreCondition_singleObj_iff}), as in the text.

\begin{center}\scriptsize\emergencystretch=3em
\begin{tabularx}{\linewidth}{>{\raggedright\arraybackslash}p{2.9cm}>{\raggedright\arraybackslash}X>{\raggedright\arraybackslash}p{2.4cm}>{\raggedright\arraybackslash}p{2.2cm}}
\toprule
statement & Lean declaration & file & status\\
\midrule
$M$-sets Boolean iff group; De Morgan iff Ore & \lean{isBooleanTopos_singleObj_iff}, \lean{isDeMorganTopos_singleObj_iff} & \lean{Monoids} & proved in Lean\\
quivers: De Morgan, Boolean (Thm~\ref{th:quivers}) & \lean{isDeMorganTopos_paths_iff}, \lean{isBooleanTopos_paths_iff} & \lean{QuiversDeMorgan} & proved in Lean (Lean's convention)\\
automata over $\Sigma$ & \lean{isDeMorganTopos_singleObj_freeMonoid_unit}, \lean{not_isDeMorganTopos_freeMonoid_two} & \lean{Monoids} & proved in Lean\\
quivers: Morita, essential points & \lean{nonempty_representations_equiv_iff}, \lean{essentialPointPathsEquiv} & \lean{QuiversMorita} & proved in Lean\\
cyclic quiver: $2^n+1$ topologies & \lean{card_grothendieckTopology} & \lean{QuiversCyclic} & proved in Lean for the count of topologies; the sheaves of each, by hand only\\
$T_n$: $n$ essential points & \lean{karoubiClassesEquivFin} & \lean{KaroubiMonoid} & proved in Lean (classes of the Karoubi envelope)\\
$M_n(\mathbb F_q)$: $n+1$ essential points & \lean{essentialPointClassesEquivFin} & \lean{EssentialPoints} & proved in Lean, any field\\
$PT_n$, $I_n$, $O_n$ essential points; all subtopos counts of \S\ref{sec:finite} & --- & --- & by hand only\\
corners: Ore transfers & \lean{forall_exists_mul_eq_mul_iff_cornerMonoid} & \lean{Corners} & proved in Lean\\
$M_n(k)$-sets, $\mathsf{Vect}_{\le n}$ & --- & --- & by hand only (the rank classes are in Lean, see above)\\
subspace lattices: De Morgan, Morita & \lean{isDeMorganTopos_functor_preorder_iff}, \lean{moritaEquivalentCat_preorder_iff} & \lean{Preorders} & proved in Lean; the rest by hand only\\
small categories: presheaf topos of a preorder iff thin (Grothendieck, 161-3, p.~26) & \lean{exists_preorder_moritaEquivalentCat_iff_isThin} & \lean{PreorderPresheaves} & proved in Lean\\
$(\mathbb Z/n,\cdot)$ counts & --- & --- & enumerated by machine (not in Lean); proof by hand only\\
compact Hausdorff $X$ & --- & --- & by hand only\\
graphs, simplicial sets, augmented, $\mathsf{Fin}$, trees & \lean{not_isDeMorganTopos_graphs}, \lean{not_isDeMorganTopos_simplicialSets}, \lean{isDeMorganTopos_augmentedSimplicialSets}, \lean{isDeMorganTopos_presheaf_fintypeCatSkeleton}, \lean{isDeMorganTopos_treeTopos} & \lean{Combinatorics} & proved in Lean; the other combinatorial rows by hand only\\
\bottomrule
\end{tabularx}
\end{center}

\subsection{How we searched}

For each statement of \S\ref{sec:machines}--\ref{sec:shapes} we searched Johnstone's \emph{Elephant}, Caramello's papers and book (including \cite{caramello-chamoun}), the papers and thesis of Hemelaer and Rogers on toposes of monoid actions \cite{hemelaer-rogers,rogers-1905,rogers-thesis}, the literature on Morita equivalence and Green's relations of monoids, the literature on synchronizing automata, the nLab, and general web and preprint searches. \q{Not found} means exactly that; several statements are elementary and may well be folklore, and we would welcome references.


\subsection{Lessons for an automatic generator}

\begin{enumerate}
\item Choosing decks is the creative step. The new outputs came from decks that we did not find paired with these arches: quivers and transformation monoids. A useful database of decks records facts of the form \q{this familiar structure is a presheaf on that shape}.
\item Unfolding is mechanical. Lemma~\ref{lem:neg} is the only rule needed to turn \q{De Morgan} into a statement about states and traps. A generator needs one such rule per invariant.
\item Rewriting is where the work goes. The raw output for automata is \q{$\B(\Sigma^*)$ satisfies \texttt{RightOreCondition}}; the readable statement is about the outgoing arrows of a vertex (Theorem~\ref{th:quivers}). The talk assigns this step to rewriting systems and language models \cite[slide 17]{caramello-mines}.
\item Filtering against the literature is necessary. Three of our candidates were found in print only at this stage (\S\ref{sec:calibration}); the number of pairs tried was not recorded, so no proportion can be given. A generator without a novelty filter reports known results as new.
\item Every output here has a direct proof: short for most, a page for the cyclic quiver, and longer for $(\mathbb Z/n,\cdot)$ and the compact Hausdorff statement, which use cited theorems (Pirashvili's classification, with a Ramsey-type factorisation; Stone and Nakano). For deeper decks a direct proof may not exist, and the formal proof of the arch in Lean is then the only guarantee, which is why the arches are formalised first.
\end{enumerate}
